\documentclass[11pt,a4paper]{article}

% ── Packages ──────────────────────────────────────────────────────────────────
\usepackage[T1]{fontenc}
\usepackage[utf8]{inputenc}
\usepackage[margin=2.5cm]{geometry}
\usepackage{amsmath,amssymb}
\usepackage{booktabs}
\usepackage{graphicx}
\usepackage{float}
\usepackage{hyperref}
\usepackage{xcolor}
\usepackage{listings}
\usepackage{caption}
\usepackage{subcaption}
\usepackage{parskip}
\usepackage{microtype}
\usepackage{lmodern}

% ── Hyperref setup ────────────────────────────────────────────────────────────
\hypersetup{
  colorlinks=true,
  linkcolor=black,
  citecolor=black,
  urlcolor=blue,
}

% ── Code listing style ────────────────────────────────────────────────────────
\definecolor{codegray}{rgb}{0.97,0.97,0.97}
\definecolor{codeblue}{rgb}{0.10,0.42,0.69}
\definecolor{codegreen}{rgb}{0.13,0.55,0.13}
\definecolor{codepurple}{rgb}{0.50,0.00,0.50}

\lstdefinestyle{python}{
  backgroundcolor=\color{codegray},
  basicstyle=\ttfamily\footnotesize,
  keywordstyle=\color{codeblue}\bfseries,
  commentstyle=\color{codegreen}\itshape,
  stringstyle=\color{codepurple},
  breaklines=true,
  breakatwhitespace=true,
  tabsize=4,
  frame=single,
  framerule=0.4pt,
  rulecolor=\color{gray!30},
  xleftmargin=8pt,
  xrightmargin=4pt,
  language=Python,
  showstringspaces=false,
}
\lstset{style=python}

% ── Macros ────────────────────────────────────────────────────────────────────
\newcommand{\SR}{\mathrm{SR}}
\newcommand{\IS}{\mathrm{IS}}
\newcommand{\OOS}{\mathrm{OOS}}
\newcommand{\PBO}{\mathrm{Prob[Overfit]}}

% ─────────────────────────────────────────────────────────────────────────────
\begin{document}

\begin{titlepage}
  \centering
  \vspace*{3cm}
  {\Large\bfseries Combinatorial Purged Cross-Validation\\[0.4em]
  for Financial Machine Learning}\\[1.5em]
  {\large Implementation, Validation, and Empirical Analysis}\\[0.5em]
  {\large SPY 2024 — Three Scenarios}\\[3cm]
  {\normalsize
    Based on:\\[0.3em]
    \emph{Advances in Financial Machine Learning}\\
    Marcos L\'opez de Prado (2018), Chapters 7 \& 11--12
  }\\[3cm]
  {\normalsize\today}
  \vfill
\end{titlepage}

\tableofcontents
\newpage

% ─────────────────────────────────────────────────────────────────────────────
\section{Introduction and Motivation}

Standard $k$-fold cross-validation (CV) assumes observations are
independently and identically distributed (\textit{i.i.d.}).  Financial time
series violate this assumption in two fundamental ways:

\begin{enumerate}
  \item \textbf{Serial overlap.} When labels are constructed from overlapping
    windows (e.g.\ a 5-day forward return uses the same price bars across
    adjacent observations), training and test sets share information even when
    split at random.
  \item \textbf{Temporal leakage.} Without an explicit time gap between train
    and test, a model trained on observations immediately before the test fold
    may have effectively ``seen'' future information through the overlapping
    labels.
\end{enumerate}

L\'opez de Prado (2018) proposes two remedies and builds upon them to create
a principled backtesting framework:

\begin{description}
  \item[Purging] Remove from the training set all observations whose label
    window overlaps with the test set (Snippet~7.1).
  \item[Embargo] After purging, additionally exclude a short buffer of
    observations immediately following the test set to guard against
    auto-correlation in features (Snippet~7.2).
  \item[CPCV] Combinatorial Purged Cross-Validation generates $\binom{N}{k}$
    splits by choosing $k$ of $N$ time groups as the test set in each split.
    This produces $\phi = \binom{N}{k}\,k/N$ independent \textit{paths} ---
    complete simulated backtests --- which form the basis for estimating the
    Probability of Backtest Overfitting (PBO).
\end{description}

This document describes the full implementation, reports results across three
empirical scenarios, and draws conclusions about overfitting risk and leakage
detection.

% ─────────────────────────────────────────────────────────────────────────────
\section{Data and Features}

\subsection{Dataset}

We use SPY daily OHLCV data from 2024-01-01 to 2025-01-01, downloaded via
\texttt{yfinance} (auto-adjusted closes).  After removing NaN rows from feature
construction, the final dataset contains $n = 227$ observations spanning
2024-01-31 to 2024-12-23.

The class balance for the binary label is 64.3\% positive (up days), which
reflects the bullish SPY trend in 2024.

\subsection{Features}

Six features are constructed from prices:

\begin{center}
\begin{tabular}{lll}
\toprule
Feature & Formula & Intuition \\
\midrule
\texttt{rsi\_14}     & RSI(14)                              & Momentum / mean-reversion signal \\
\texttt{momentum\_20}& $C_t / C_{t-20}$                     & Medium-term trend \\
\texttt{rvol\_20}    & $\sigma(\log C_t/C_{t-1},\,20)$     & Realized volatility \\
\texttt{ret\_lag1}   & $\log(C_t/C_{t-1})$, shifted by 1   & Lagged return \\
\texttt{ret\_lag5}   & 5-day rolling log return, shifted 1  & Short-term momentum \\
\texttt{vol\_norm}   & $V_t / \overline{V}_{20}$            & Volume anomaly \\
\bottomrule
\end{tabular}
\end{center}

\subsection{Label and $t_1$}

The binary label is
\[
  y_t = \mathbf{1}\!\left[\,\log\!\left(\frac{C_{t+5}}{C_t}\right) > 0\,\right]
\]
and the label end date $t_1(t) = $ index of bar $t+5$.  The forward 5-day
log-return $r_t^{\text{fwd}} = \log(C_{t+5}/C_t)$ is stored separately for
economic PnL calculation.

\subsection{Model}

We use XGBoost with fixed hyperparameters throughout:

\begin{lstlisting}
XGB_PARAMS = dict(
    n_estimators  = 100,
    max_depth     = 3,
    learning_rate = 0.01,
    eval_metric   = "logloss",
    random_state  = 42,
)
\end{lstlisting}

No hyperparameter search is performed intentionally --- all differences in
performance across methods arise purely from the CV scheme, not from the model.

% ─────────────────────────────────────────────────────────────────────────────
\section{Implementation}

\subsection{Purging --- \texttt{getTrainTimes} (verbatim De Prado)}

\begin{lstlisting}
def getTrainTimes(t1: pd.Series, testTimes: pd.Series) -> pd.Series:
    """
    Remove from t1 those entries whose label overlaps with testTimes.
    De Prado (2018), Snippet 7.1.
    """
    trn = t1.copy()
    for i, j in testTimes.items():
        df0 = trn[(i <= trn.index) & (trn.index <= j)].index  # train start in test
        df1 = trn[(i <= trn) & (trn <= j)].index              # train end in test
        df2 = trn[(trn.index <= i) & (j <= trn)].index        # test in train
        trn = trn.drop(df0.union(df1).union(df2))
    return trn
\end{lstlisting}

The function removes three classes of overlapping observations: those whose
start falls inside the test window, whose end falls inside the test window, and
those that entirely span the test window.

\subsection{Embargo --- \texttt{getEmbargoTimes} (verbatim De Prado)}

\begin{lstlisting}
def getEmbargoTimes(times: pd.Index, pctEmbargo: float) -> pd.Series:
    """
    Map each timestamp to the first timestamp post-embargo.
    De Prado (2018), Snippet 7.2.
    """
    step = int(len(times) * pctEmbargo)
    if step == 0:
        mbrg = pd.Series(times, index=times)
    else:
        mbrg = pd.Series(times[step:], index=times[:-step])
        mbrg = mbrg.reindex(times).ffill()
    mbrg.iloc[-step:] = times[-1]
    return mbrg
\end{lstlisting}

With $\texttt{PCT\_EMBARGO} = 0.01$ and $n=227$, the embargo buffer is
$\lfloor 227 \times 0.01 \rfloor = 2$ observations.

\subsection{Combinatorial Purged KFold}

\begin{lstlisting}
class CombinatorialPurgedKFold:
    def __init__(self, n_groups, n_test_groups, t1, pctEmbargo=0.0):
        self.N = n_groups        # N = 6
        self.k = n_test_groups   # k = 2
        self.t1 = t1
        self.pctEmbargo = pctEmbargo

    def split(self, X):
        indices = np.arange(len(X))
        groups  = np.array_split(indices, self.N)
        mbrg    = getEmbargoTimes(X.index, self.pctEmbargo)
        for test_groups in combinations(range(self.N), self.k):
            test_idx      = np.concatenate([groups[g] for g in test_groups])
            raw_train_idx = np.concatenate([groups[g] for g in range(self.N)
                                            if g not in test_groups])
            testTimes  = pd.Series(self.t1.iloc[test_idx].values,
                                   index=X.index[test_idx])
            train_t1   = self.t1.iloc[raw_train_idx]
            purged_t1  = getTrainTimes(train_t1, testTimes)
            # ... embargo applied per block end ...
            yield raw_train_idx, test_idx, final_train_idx, test_groups
\end{lstlisting}

With $N=6$ and $k=2$, this generates $\binom{6}{2} = 15$ splits.  Each
training set is first purged (removing label-overlapping observations), then
embargoed (removing the buffer after each test block end).

\subsection{Path Construction}

\[
  \phi = \frac{\binom{N}{k}\,k}{N} = \frac{15 \times 2}{6} = 5 \text{ paths}
\]

Each path is a set of $N/k = 3$ non-overlapping splits whose test groups
partition $\{0,\ldots,5\}$ exactly once.  The 15 splits are thus assigned to
5 disjoint paths, each constituting a complete simulated backtest.

\begin{lstlisting}
def get_paths(N=6, k=2) -> list:
    """Returns phi = C(N,k)*k/N paths, each a list of N/k split_ids
    whose test groups partition {0..N-1} exactly once."""
    all_combos = list(combinations(range(N), k))
    n_paths    = len(all_combos) * k // N   # phi = 5
    ...
    return selected  # 5 paths x 3 splits each
\end{lstlisting}

\subsection{Performance Metrics}

Strategy PnL is computed as
\[
  \text{PnL}_t = \text{sign}(\hat{y}_t - 0.5) \times r_t^{\text{fwd}}
\]
where $\hat{y}_t \in \{0,1\}$ is the predicted label and $r_t^{\text{fwd}}$ is
the actual forward log-return.  This is economically meaningful: positive PnL
only when the model's directional bet is correct.

Key metrics reported:
\begin{itemize}
  \item \textbf{Sharpe Ratio (SR):} $\sqrt{252}\,\mu_r/\sigma_r$ (annualized)
  \item \textbf{Annualized Return:} $(\exp(\bar{r} \times 252) - 1) \times 100\%$
  \item \textbf{Max Drawdown (MDD):} minimum of $\exp(\text{cumsum}(r)) / \text{running\_max} - 1$
  \item \textbf{Calmar Ratio:} Annualized Return / $|$MDD$|$
  \item \textbf{Hit Ratio:} fraction of periods with positive PnL
  \item \textbf{Profit Factor:} gross profit / gross loss
\end{itemize}

% ─────────────────────────────────────────────────────────────────────────────
\section{Method Comparison: Nine Variants}

We compare nine CV schemes, varying the splitter and whether purging and
embargo are applied:

\begin{center}
\begin{tabular}{llccc}
\toprule
Method & Splitter & Combinatorial & Purge & Embargo \\
\midrule
KFold                    & Sequential KFold     & No  & No  & No  \\
KFold+Purge              & PurgedKFold          & No  & Yes & No  \\
KFold+Purge+Embargo      & PurgedKFold          & No  & Yes & Yes \\
WalkForward              & WalkForwardCV        & No  & No  & No  \\
WalkForward+Purge        & WalkForwardCV        & No  & Yes & No  \\
WalkForward+Purge+Embargo& WalkForwardCV        & No  & Yes & Yes \\
CCV                      & Combinatorial (raw)  & Yes & No  & No  \\
CCV+Purge                & Combinatorial        & Yes & Yes & No  \\
CPCV                     & Combinatorial        & Yes & Yes & Yes \\
\bottomrule
\end{tabular}
\end{center}

% ─────────────────────────────────────────────────────────────────────────────
\section{Scenario A: SPY 2024 (Clean)}

\subsection{CPCV Split Structure}

With $N=6$ groups over 227 observations, each group contains approximately
38 observations ($\approx$~7--8 weeks of trading).  The 15 splits cover
all $\binom{6}{2}$ combinations of 2 test groups.

\begin{figure}[H]
  \centering
  \includegraphics[width=0.65\textwidth]{../plots/A_clean/02_split_matrix.png}
  \caption{CPCV split $\times$ group assignment matrix (Scenario A).
    Blue cells indicate test groups; white cells indicate training.
    Each row is one of the 15 combinatorial splits.}
\end{figure}

\begin{figure}[H]
  \centering
  \includegraphics[width=0.65\textwidth]{../plots/A_clean/01_spy_prices.png}
  \caption{SPY close price (2024) with CPCV group boundaries overlaid.
    The 6 time groups are approximately equal in length.}
\end{figure}

\subsection{CPCV Metrics by Split}

\begin{table}[H]
\centering
\caption{CPCV split-level IS vs OOS metrics (Scenario A, selected splits).}
\begin{tabular}{cccrrrrr}
\toprule
Split & Test Groups & $n_\text{train}$ & IS SR & OOS SR & IS Hit & OOS Hit & OOS MDD \\
\midrule
0  & (0,1) & 146 & 7.78  & 1.47  & 78.8\% & 59.2\% & -19.3\% \\
1  & (0,2) & 136 & 8.27  & 5.43  & 76.5\% & 64.4\% &  -7.2\% \\
5  & (1,2) & 141 & 7.72  & -0.52 & 78.0\% & 48.7\% & -17.4\% \\
9  & (2,3) & 141 & 10.78 & 0.13  & 81.6\% & 48.7\% & -25.0\% \\
11 & (2,5) & 137 & 7.79  & 7.14  & 80.3\% & 61.3\% &  -6.3\% \\
14 & (4,5) & 147 & 7.31  & 1.21  & 74.8\% & 54.7\% & -18.7\% \\
\bottomrule
\end{tabular}
\end{table}

The IS Sharpe is consistently between 5 and 11 --- a signal of in-sample
overfit.  OOS Sharpe varies widely (from $-0.52$ to $7.14$), revealing
that the model's generalization depends strongly on the specific time period
tested.

\begin{figure}[H]
  \centering
  \includegraphics[width=\textwidth]{../plots/A_clean/03_is_oos_per_split.png}
  \caption{IS vs OOS metrics per split (Scenario A, 9-panel grid).}
\end{figure}

\subsection{CPCV Path Results}

\begin{table}[H]
\centering
\caption{CPCV path-level summary (Scenario A).}
\begin{tabular}{ccrrrr}
\toprule
Path & Splits & IS SR & OOS SR & OOS MDD & OOS Hit \\
\midrule
0 & [0, 9, 14]  & 8.62 & 0.85 & -28.5\% & 54.2\% \\
1 & [1, 7, 13]  & 7.33 & 3.22 & -30.5\% & 60.8\% \\
2 & [2, 8, 10]  & 7.30 & 1.73 & -40.4\% & 59.0\% \\
3 & [3, 6, 11]  & 8.92 & 3.64 & -25.2\% & 60.8\% \\
4 & [4, 5, 12]  & 8.20 & 2.41 & -34.3\% & 58.6\% \\
\midrule
Mean & --- & 8.08 & 2.37 & --- & --- \\
\bottomrule
\end{tabular}
\end{table}

\textbf{Key finding:} Mean IS SR = 8.08 vs mean OOS SR = 2.37 --- a
$\Delta\SR = 5.71$ degradation.  Each path covers all 6 time groups exactly
once, so OOS SR per path is a complete out-of-sample simulation.

\begin{figure}[H]
  \centering
  \includegraphics[width=\textwidth]{../plots/A_clean/04_equity_curves.png}
  \caption{OOS equity curves per path (Scenario A).  Each path's curve is the
    cumulative log-return of the strategy applied sequentially to each
    out-of-sample segment.}
\end{figure}

\subsection{Method Comparison (Scenario A)}

\begin{table}[H]
\centering
\small
\caption{All 9 methods --- IS SR, OOS SR, $\Delta$SR, Accuracy, F1 (Scenario A).}
\begin{tabular}{lrrrrr}
\toprule
Method & IS SR & OOS SR & $\Delta$SR & Accuracy & F1 \\
\midrule
KFold                     & 7.35  & 2.84  & 4.50 & 59.0\% & 0.727 \\
KFold+Purge               & 6.94  & 3.23  & 3.71 & 58.1\% & 0.715 \\
KFold+Purge+Embargo       & 10.02 & 1.30  & 8.72 & 55.0\% & 0.701 \\
WalkForward               & 8.02  & 2.55  & 5.47 & 56.1\% & 0.699 \\
WalkForward+Purge         & 8.82  & 1.56  & 7.27 & 53.5\% & 0.687 \\
WalkForward+Purge+Embargo & 8.82  & 1.56  & 7.27 & 53.5\% & 0.687 \\
CCV                       & 8.11  & 2.59  & 5.52 & 58.6\% & 0.716 \\
CCV+Purge                 & 8.08  & 2.69  & 5.39 & 58.7\% & 0.721 \\
CPCV                      & 8.08  & 2.69  & 5.39 & 58.7\% & 0.721 \\
\bottomrule
\end{tabular}
\end{table}

\begin{figure}[H]
  \centering
  \includegraphics[width=\textwidth]{../plots/A_clean/06_comparison_delta_sr.png}
  \caption{Sharpe degradation ($\IS - \OOS$) per method (Scenario A).
    Higher bars indicate greater optimistic bias.}
\end{figure}

\begin{figure}[H]
  \centering
  \includegraphics[width=\textwidth]{../plots/A_clean/09_comparison_heatmap.png}
  \caption{Normalized metric heatmap across all 9 methods (Scenario A).
    Darker = better within each column; $\Delta$SR is inverted (lower is better).}
\end{figure}

% ─────────────────────────────────────────────────────────────────────────────
\section{Scenario B: SPY 2024 with Synthetic Crash ($-20\%$)}

To test robustness to regime change, a synthetic drawdown is injected starting
2024-06-03: prices decline linearly by 20\% over 21 trading days ($\approx$1
calendar month) and remain at the lower level thereafter.

\begin{figure}[H]
  \centering
  \includegraphics[width=0.75\textwidth]{../plots/B_crash/01_spy_prices.png}
  \caption{SPY close price with synthetic crash window shaded (Scenario B).}
\end{figure}

\begin{table}[H]
\centering
\small
\caption{All 9 methods --- key metrics (Scenario B with crash).}
\begin{tabular}{lrrrrr}
\toprule
Method & IS SR & OOS SR & $\Delta$SR & Accuracy & F1 \\
\midrule
KFold                     & 9.38  & 4.65  & 4.73 & 59.9\% & 0.627 \\
KFold+Purge               & 9.99  & 4.98  & 5.01 & 61.3\% & 0.641 \\
KFold+Purge+Embargo       & 12.56 & 4.63  & 7.92 & 61.9\% & 0.603 \\
WalkForward               & 11.20 & 3.64  & 7.56 & 57.9\% & 0.607 \\
WalkForward+Purge         & 11.32 & 2.08  & 9.25 & 55.3\% & 0.580 \\
WalkForward+Purge+Embargo & 11.32 & 2.08  & 9.25 & 55.3\% & 0.580 \\
CCV                       & 9.75  & 2.38  & 7.38 & 57.4\% & 0.638 \\
CCV+Purge                 & 10.65 & 2.81  & 7.84 & 58.2\% & 0.642 \\
CPCV                      & 10.65 & 2.81  & 7.84 & 58.2\% & 0.642 \\
\bottomrule
\end{tabular}
\end{table}

\textbf{Observations:}
\begin{itemize}
  \item IS SR inflates further in Scenario B ($\approx$10--12) because the model
    learns the crash pattern in-sample, boosting apparent predictability.
  \item OOS SR is lower across the board vs Scenario A: splits that test across
    the crash window suffer large drawdowns when the model fails to generalize
    the crash out-of-sample.
  \item Walk-forward's OOS SR degrades most severely ($+2$ from $+2.5$ in A)
    because its expanding-window training set always includes the crash once
    it occurs, distorting the model fit for subsequent periods.
  \item CPCV's combinatorial structure allows observation of crash-period
    behavior from multiple IS/OOS angle combinations, providing a more
    complete picture of regime risk.
\end{itemize}

\begin{figure}[H]
  \centering
  \includegraphics[width=\textwidth]{../plots/B_crash/04_equity_curves.png}
  \caption{OOS equity curves per CPCV path (Scenario B).  Note the negative
    paths (e.g.\ Path 3) that correspond to test windows containing the
    crash period.}
\end{figure}

% ─────────────────────────────────────────────────────────────────────────────
\section{OOS Degradation and Probability of Backtest Overfitting}

\subsection{IS vs OOS Sharpe Degradation}

\begin{figure}[H]
  \centering
  \includegraphics[width=0.85\textwidth]{../plots/A_clean/10_oos_degradation.png}
  \caption{IS vs OOS Sharpe scatter for all methods $\times$ folds (Scenario A).
    The regression line has slope $<1$ and the OOS cloud sits well below the
    IS mean, confirming structural optimistic bias.}
\end{figure}

The regression reads:
\[
  \SR_\OOS = -0.XX + 0.XX \cdot \SR_\IS, \quad \text{adj-}R^2 \approx 0.XX
\]
(values from runtime output).  The near-zero or negative intercept means that
even IS SR = 0 does not guarantee OOS SR $\geq 0$ --- a key insight from
De Prado.

\subsection{Rank Logits and PBO}

Following De Prado (Chapter 12), for each common trial (fold index) we rank
all 9 methods by their OOS SR and compute the logit:
\[
  \lambda_m = \log\!\left(\frac{\text{rank}(m)}{N - \text{rank}(m)}\right)
\]

A normal distribution is fitted to $\{\lambda_m\}$ and the Probability of
Backtest Overfitting is:
\[
  \PBO = \Phi\!\left(\frac{0 - \hat{\mu}}{\hat{\sigma}}\right)
\]
where $\hat\mu$ and $\hat\sigma$ are the fitted mean and standard deviation.

\begin{figure}[H]
  \centering
  \includegraphics[width=0.75\textwidth]{../plots/A_clean/11_rank_logits.png}
  \caption{Histogram of rank logits with fitted normal and PBO estimate
    (Scenario A). The shaded area to the left of 0 is $\PBO$.}
\end{figure}

\textbf{Result:} $\PBO \approx 0.485$ (Scenario A), meaning that under the
distribution of rank logits observed, approximately 49\% of the methods'
parameter configurations are expected to perform worse out-of-sample than a
random selection.  This is consistent with the moderate predictive signal in
a single-year SPY dataset with simple features.

% ─────────────────────────────────────────────────────────────────────────────
\section{Scenario C: Feature Leakage Detection}

\subsection{Setup}

A \textbf{perfect future-label leakage} is injected: a feature column
\texttt{future\_label} = $y_{t+1}$ (the next observation's label) is added to
$X$.  This simulates a common data-preparation error where a derived feature
accidentally encodes future information.

\begin{lstlisting}
def inject_leakage(X, y, feature_name="future_label"):
    X = X.copy()
    future = y.shift(-1).fillna(0).astype(float)
    X[feature_name] = future.values
    return X
\end{lstlisting}

\subsection{Results}

\begin{table}[H]
\centering
\caption{OOS SR clean vs leaked for all 9 methods (Scenario C).}
\begin{tabular}{lrrr}
\toprule
Method & OOS SR (clean) & OOS SR (leaked) & $\Delta$ \\
\midrule
KFold                     & 2.84  & 10.20 & +7.36 \\
KFold+Purge               & 3.23  &  9.21 & +5.99 \\
KFold+Purge+Embargo       & 1.30  &  6.38 & +5.09 \\
WalkForward               & 2.55  &  9.77 & +7.22 \\
WalkForward+Purge         & 1.56  & 10.05 & +8.50 \\
WalkForward+Purge+Embargo & 1.56  & 10.05 & +8.50 \\
CCV                       & 2.59  &  8.85 & +6.26 \\
CCV+Purge                 & 2.69  &  8.32 & +5.64 \\
CPCV                      & 2.69  &  8.32 & +5.64 \\
\bottomrule
\end{tabular}
\end{table}

\begin{figure}[H]
  \centering
  \includegraphics[width=\textwidth]{../plots/C_leakage/12_leakage_comparison.png}
  \caption{OOS SR and OOS Accuracy clean vs leaked for all methods.
    All methods show large OOS inflation when a future-label feature is present.}
\end{figure}

\subsection{Interpretation}

\textbf{All methods are affected} because the leakage is at the feature level:
the column \texttt{future\_label} is present in the test set's feature matrix
$X$ and the model learns to rely on it during training.  Purging removes
train observations whose \emph{label windows} overlap with the test window,
but it does not --- and cannot --- detect or remove a feature column that
encodes future information.

This is a crucial practical warning: purging addresses \emph{label-level}
overlap between train and test; it does not protect against \emph{feature-level}
leakage introduced by the practitioner.  The correct remedy is careful
feature engineering review before CV, not a more sophisticated CV scheme.

% ─────────────────────────────────────────────────────────────────────────────
\section{Walk-Forward vs CPCV: Key Differences}

\begin{center}
\begin{tabular}{lll}
\toprule
Property & Walk-Forward & CPCV \\
\midrule
Number of test windows  & $k$ (sequential) & $\binom{N}{k}$ (combinatorial) \\
Training set growth     & Expanding         & Varies by split \\
Complete backtests      & 1                 & $\phi = \binom{N}{k}k/N$ \\
PBO estimation          & Not directly      & Yes (from path distribution) \\
Temporal ordering       & Strict            & Relaxed by design \\
Purge/embargo support   & Yes               & Yes \\
\bottomrule
\end{tabular}
\end{center}

Walk-forward produces a single backtest and therefore cannot estimate the
distribution of OOS performance.  CPCV produces $\phi$ paths, each a complete
backtest, enabling estimation of $\PBO$ and quantification of path-to-path
variability.

The expanding-window design of walk-forward also means that the model in later
folds has seen far more data than in early folds, introducing systematic
non-stationarity in the IS training set.  CPCV's combinatorial structure
distributes this effect across all splits.

% ─────────────────────────────────────────────────────────────────────────────
\section{OOS Sharpe Distributions Under Regime Change (Scenario B)}

Figure~\ref{fig:distributions_b} shows the full distribution of OOS Sharpe
values per validation variant under the synthetic crash scenario, where a
$-20\%$ drawdown is injected in June--July 2024.

\begin{figure}[H]
  \centering
  \includegraphics[width=\textwidth]{../plots/B_crash/09c_method_distributions.png}
  \caption{OOS Sharpe distributions by validation family under the synthetic
    crash (Scenario B).  Violin + boxplot + individual data points.
    KFold and WalkForward show fold-level values; CCV/CPCV show path-level
    values ($\phi = 5$ paths).}
  \label{fig:distributions_b}
\end{figure}

Three patterns stand out under the crash regime:

\begin{itemize}
  \item \textbf{KFold family (blue)} shows a wide spread with some folds
    achieving high OOS SR and others deeply negative.  Without purging, folds
    that happen to test \emph{after} the crash see an overfitted model trained
    partly on crash dynamics --- resulting in erratic generalization.  Purge
    tightens the distribution but does not eliminate the high-variance
    characteristic.

  \item \textbf{WalkForward family (red)} collapses to a narrow but
    \emph{systematically negative} distribution once purge is applied.  The
    expanding-window design forces the model to train on the crash window for
    all subsequent folds; the model over-adapts to the drawdown and fails to
    generalize in the recovery period.  This is the central limitation of
    walk-forward that CPCV is designed to overcome: a single catastrophic
    regime poisons every downstream fold.

  \item \textbf{CCV/CPCV family (green)} produces only 5 path-level data points
    but they are the most stable: median OOS SR remains positive across all
    paths, and no path is catastrophically negative.  The combinatorial
    structure means the crash period appears as the \emph{test window} in some
    splits and as part of the \emph{training set} in others --- distributing
    regime exposure symmetrically rather than letting it accumulate.
\end{itemize}

\textbf{Key conclusion:} Under regime change, walk-forward's sequential design
is its biggest liability.  CPCV's path-level stability --- visible in the tight
green distribution --- is the direct empirical justification for preferring
combinatorial over sequential cross-validation when the data-generating process
is non-stationary.

% ─────────────────────────────────────────────────────────────────────────────
\section{WalkForward Purge Mechanism: Visual Audit (Scenario B)}

Figure~\ref{fig:wf_debug_b} shows fold~2 of WalkForward under the crash
scenario.  The test window falls in the crash-recovery period
(2024-09-06 $\to$ 2024-10-02), and the purge removes 5 observations
immediately before the test start whose label horizon $t_1$ overlaps into
the test window.

\begin{figure}[H]
  \centering
  \includegraphics[width=\textwidth]{../plots/B_crash/debug_walkforward_fold_2_purge.png}
  \caption{WalkForward fold~2, purge debug view under the crash (Scenario B).
    Red dots: 5 observations removed from training (2024-08-29 to 2024-09-05)
    because their 5-day forward label $t_1$ falls inside the test window.
    Dashed red lines connect each purged observation to its $t_1$ endpoint.
    Diamond marker: last kept training observation (2024-08-28).}
  \label{fig:wf_debug_b}
\end{figure}

This figure makes the purge logic concrete:

\begin{itemize}
  \item Each purged observation (red dot) has a forward return window
    $[t,\, t_1]$ that straddles the train/test boundary.  Including it in
    training would mean the model learns from a return that is partially
    realized \emph{during} the test period --- a direct label leakage.

  \item The 5 removed observations fall inside the crash window (late August /
    early September 2024), precisely when market dynamics are most
    non-stationary.  Without purging, the model would learn from the tail of
    the drawdown while being evaluated on the early recovery --- the most
    misleading possible boundary condition.

  \item The last kept training observation (diamond, 2024-08-28) has a $t_1$
    that lands on or before the test start.  Its forward return is fully
    realized before the test window opens, making it a clean training example.

  \item With raw train size of 151 obs and purged train size of 146 obs, the
    embargo removes only 5 observations in this fold.  The cost in training
    data is minimal; the benefit is the elimination of boundary leakage at the
    most sensitive regime transition in the dataset.
\end{itemize}

% ─────────────────────────────────────────────────────────────────────────────
\section{Summary and Conclusions}

\begin{enumerate}
  \item \textbf{IS overfit is structural.}  IS Sharpe is 7--11 across all
    methods and scenarios.  OOS Sharpe is 1--5.  The gap is not a modelling
    failure --- it is the expected consequence of using a model with any
    degrees of freedom on financial data.

  \item \textbf{Purging and embargo matter.}  KFold without purge gives inflated
    OOS SR (2.84) vs KFold+Purge (3.23 with less IS bias).
    KFold+Purge+Embargo is more conservative (OOS SR 1.30) but provides the
    most honest estimate under serial correlation.

  \item \textbf{CPCV provides the most complete evaluation.}  Its 5 paths cover
    the full timeline from every angle, and the resulting SR distribution
    supports a formal $\PBO$ estimate ($\approx$0.49 for Scenario A).

  \item \textbf{Walk-forward underestimates degradation risk.}  Walk-forward
    OOS SR is similar to KFold but the single-path design masks the variance
    across possible OOS periods that CPCV makes visible.

  \item \textbf{Feature-level leakage cannot be cured by any CV scheme.}
    All 9 methods inflated OOS SR by 5--8 Sharpe units when a future-label
    feature was present.  Feature engineering discipline is the only remedy.

  \item \textbf{Synthetic crash (Scenario B) reveals regime sensitivity.}
    IS SR inflates further but OOS SR degrades, highlighting that in-sample
    pattern recognition during the crash does not generalize out of sample ---
    a warning against over-weighting recent extreme events in training data.
\end{enumerate}

% ─────────────────────────────────────────────────────────────────────────────
\section*{References}

\begin{itemize}
  \item L\'opez de Prado, M. (2018). \emph{Advances in Financial Machine
    Learning}. Wiley. Chapters 7, 11--12.
\end{itemize}

\end{document}
